\documentclass{jsarticle}
\usepackage{amsmath}
\usepackage{amssymb}
\begin{document}
\title{}
\author{}
\date{}
\maketitle

  自明な零点は実軸${1 \over 2} $で、0は勾配率$x=0$で
  等加速度、等速度、接平面$M=0,1,..,x$を値として、
  特殊相対性理論、一般相対性理論に 崩壊定理、不動点定理で 特異点0、
  ホモロジー、コホモロジー、  オイラー数が2であり、
  自明な零点でentropy値を一意に原子と同じくとる。
  相加相乗平均は、実部と虚部${{1 \over 2}+iy}$に値として、
  三角関数と逆三角関数をオイラーの公式でガンマ関数とベータ関数
  と同じく、逆関数で虚数と実数の入れ換えをして、正規部分群
  で極限値をとると、可積分系の超弦理論の場の理論となる。
  世界面を曲平面として、世界樹が宇宙と異次元空間を原子の構造と同じくし、
  D-braneとAnti-D-braneを原子同士を集めると同じくassemble-D-braneを
  構築して、血脈と同じ仕組みで、宇宙と異次元空間の組合せ多様体
  が、全部で6種類ある。2種類ずつの次元が2世界あり、1種類の次元で
  ザイフェルト多様体を構築し、このザイフェルト多様体が双対分割で、
  2種類ある。超対称性素粒子が全部で12種類あるのが、この次元の
  組合せ多様体の世界樹と同じとしてある。




\end{document}
